\documentclass[12pt, a4paper, leqno]{article}

\usepackage[T1]{fontenc}
\usepackage[utf8]{inputenc}

% Font: Times New Roman (serif, as required)
\usepackage{times}  % or \usepackage{mathptmx} for math support

\usepackage{float, afterpage, rotating, graphicx}
\usepackage{epstopdf}
\usepackage{longtable, booktabs, tabularx}
\usepackage{fancyvrb, moreverb, relsize}
\usepackage{eurosym, calc}
\usepackage{amsmath, amssymb, amsfonts, amsthm, bm}
\usepackage{caption}
\usepackage{mdwlist}
\usepackage{xfrac}
\usepackage{setspace}
\usepackage[dvipsnames]{xcolor}
\usepackage{subcaption}
\usepackage{minibox}

% Margins: 3cm left, 2cm right, top and bottom
\usepackage[
    left=3cm,
    right=2cm,
    top=2cm,
    bottom=2cm
]{geometry}

% Line spacing: 1.5 as required
\usepackage{setspace}
\setstretch{1.5}

% Footnote font size (size 8 or 10, single/basic spacing)
\usepackage{footmisc}
\renewcommand{\footnotelayout}{\fontsize{10}{12}\selectfont\singlespacing}

% Bibliography
\usepackage[
    natbib=true,
    bibencoding=inputenc,
    bibstyle=authoryear-ibid,
    citestyle=authoryear-comp,
    maxcitenames=3,
    maxbibnames=10,
    useprefix=false,
    sortcites=true,
    backend=biber
]{biblatex}
\AtBeginDocument{\toggletrue{blx@useprefix}}
\AtBeginBibliography{\togglefalse{blx@useprefix}}
\setlength{\bibitemsep}{1.5ex}
\addbibresource{refs.bib}

% Hyperlinks
\usepackage[unicode=true]{hyperref}
\hypersetup{
    colorlinks=true,
    linkcolor=black,
    anchorcolor=black,
    citecolor=NavyBlue,
    filecolor=black,
    menucolor=black,
    runcolor=black,
    urlcolor=NavyBlue
}

% Pagination: Arabic numerals, starting at 1 after TOC
\usepackage{fancyhdr}
\pagestyle{fancy}
\fancyhf{}
\fancyfoot[C]{\thepage}
\renewcommand{\headrulewidth}{0pt}

\widowpenalty=10000
\clubpenalty=10000
\setlength{\parskip}{1ex}
\setlength{\parindent}{0ex}

% Full justification (LaTeX default, but explicit)
\usepackage{ragged2e}

\begin{document}

% -------------------------------------------------------
% COVER SHEET (unnumbered, no page number shown)
% -------------------------------------------------------
\begin{titlepage}
    \centering
    \vspace*{3cm}
    {\large Lifecycle Consumption under Budget Constraints: A Replication of 'Estimating Discount Functions with Consumption Choices over the Lifecycle' \par}
    \vspace{2cm}
    {\normalsize
    Master Thesis Presented to the \\
    Department of Economics at the \\
    Rheinische Friedrich-Wilhelms-Universität Bonn \\[1em]
    In Partial Fulfillment of the Requirements for the Degree of \\
    Master of Science (M.Sc.) \par}
    \vspace{2cm}
    {\normalsize Supervisor: Hans-Martin von Gaudecker \par}
    \vspace{1cm}
    {\normalsize Submitted in May, 2026 by: \\[0.5em]
    Brenda Torres Blomer \\
    Matriculation Number: [Number] \par}
\end{titlepage}

% -------------------------------------------------------
% TABLE OF CONTENTS (unnumbered pages, no Arabic counter)
% -------------------------------------------------------
\pagenumbering{gobble}   % suppress page numbers on TOC
\tableofcontents
\newpage

% Optional lists (only if needed):
% \listoffigures
% \listoftables
% \newpage

% -------------------------------------------------------
% MAIN BODY: Arabic pagination starts at 1
% -------------------------------------------------------
\pagenumbering{arabic}
\setcounter{page}{1}

\section{Introduction}
\label{sec:introduction}

Behavioral economics studies how individual decisions can systematically deviate from the predictions of standard rational models. One important example is present bias, where individuals place disproportionate weight on immediate consumption relative to future consumption. In practice, this can lead individuals to mispredict their future behavior. For example, someone may choose a cheaper mobile data plan believing that their future usage will remain low, only to later consume more data than expected and pay higher additional costs.

To capture this type of behavior, economists developed models that incorporate time-inconsistent preferences. One of the most influential frameworks is the $\beta$-$\delta$ model of quasi-hyperbolic discounting, introduced by \citet{laibson1997golden}. In this framework, individuals may display present bias by valuing immediate utility more heavily than future utility. Experimental and empirical evidence has shown that this type of behavior can help explain several consumption and saving patterns observed in the data.

The paper \textit{Estimating Discount Functions with Consumption Choices over the Lifecycle} by \citet{laibson1997golden} studies whether present-biased preferences can explain observed household patterns of debt accumulation and illiquid wealth holdings over the lifecycle. The authors develop and estimate a lifecycle model in which households may behave as naive present-biased agents and compare its performance with a standard exponential discounting model. Their results suggest that the present-biased specification provides a better fit to the empirical lifecycle moments.

Lifecycle models of this type are computationally complex and rely heavily on numerical methods, discretization choices, and simulation procedures. For this reason, replication plays an important role not only for verifying empirical findings but also for understanding how sensitive results may be to alternative computational implementations. Even when authors provide detailed replication packages, reproducing the exact results of structural lifecycle models can remain challenging due to differences in software architecture, optimization procedures, stochastic simulations, and numerical approximations.

This thesis conducts a computational replication of \citet{laibson2024estimating} using PyLCM from \citet{pylcm2024}, a Python package designed for modular lifecycle modeling. The objective is not to modify the underlying economic model, but rather to evaluate whether the main lifecycle patterns and simulated moments of the original paper can be reproduced under an alternative computational framework. In particular, this thesis focuses on how differences in grid construction, stochastic processes, constraints, and simulation architecture affect the implementation and resulting lifecycle dynamics.

The thesis proceeds as follows. Section \ref{sec:structure} summarizes the original paper and the parts of the model most relevant for the replication. Section \ref{sec:replication} describes the implementation in PyLCM, the computational structure of the model, and the numerical challenges encountered during the replication process. Section \ref{sec:results} presents the simulation results and compares them with the moments reported in the original paper. Finally, the last section discusses the main conclusions, limitations, and possible directions for future research.

%And, for the final section I add my extension that is just a start for later research that could be use if some technical problems are solved ...ver si pongo esto ultimo.


\section{Overview of the Original Study}
\label{sec:structure}

\subsection{Research Objective}

The paper I replicate is ``Estimating Discount Functions with Consumption Choices over the Lifecycle'' by \citet{laibson2024estimating}. The main objective of the paper is to study whether present-biased preferences can explain observed patterns of borrowing, saving, and wealth accumulation over the life cycle. In particular, the authors focus on households that simultaneously accumulate illiquid wealth while carrying high-interest credit card debt, a behavior that could be difficult to explain with standard lifecycle models under exponential discounting.

To study this question, the paper develops a structural lifecycle model in which households make consumption, borrowing, and saving decisions throughout their lives under income uncertainty. The model generates simulated moments related to debt holdings, wealth accumulation, and credit card use, which are then compared with empirical moments from United States (U.S.) household surveys in order to estimate households' preference parameters.

The paper contributes to the literature by incorporating present-biased preferences into a lifecycle framework, allowing the model to capture behavioral patterns that standard models struggle to explain.

\subsection{Lifecycle Model}

The paper follows the lifecycle model developed by \citet{angeletos2001hyperbolic}. In the model, households enter the economy at age 20 and live up to age 90. Throughout the lifecycle, households face uncertainty in income, survival probabilities, and demographic composition, including changes in the number of children and dependent adults. Income follows a stochastic process with both persistent and transitory shocks, generating heterogeneous income paths across households.

Households can save and borrow using two different types of assets. The first is a liquid asset, denoted by $X_t$, which can take both positive and negative values. Positive liquid wealth earns a saving interest rate, while negative liquid wealth represents credit card debt and accrues a higher borrowing interest rate. The amount households can borrow depends on a credit limit that varies with age. The second asset is an illiquid asset, denoted by $Z_t$. Households are not allowed to hold negative illiquid wealth. Unlike liquid wealth, this asset does not generate direct capital gains, but it provides an annual return to consumption, $\gamma = 0.05$, proportional to the amount held. Selling illiquid wealth is costly, with the liquidity cost depending on age, $\kappa_t = \frac{0.5}{1+e^{(t-50)/10}}$.

Households choose how much to allocate to liquid and illiquid wealth each period. Let $I_t^X$ denote the investment decision for liquid wealth and $I_t^Z$ the investment decision for illiquid wealth. Negative values of $I_t^Z$ represent sales of illiquid assets and imply paying the liquidity cost. The resulting consumption flow is given by

\begin{equation}
C_t = Y_t - I_t^X - I_t^Z + \kappa_t \min(I_t^Z,0).
\end{equation}

Households derive utility from consumption and from holding illiquid wealth. Preferences follow a Constant Relative Risk Aversion (CRRA) utility specification with risk aversion coefficient $\rho$:

\begin{equation}
u(C_t)=n_t\cdot\frac{\left(\frac{C_t+\gamma Z_t}{n_t}\right)^{1-\rho}-1}{1-\rho}.
\end{equation}

The term $\gamma Z_t$ captures the return on consumption generated by holding illiquid wealth. In addition, households derive utility from leaving bequests at the end of life.

The model allows for both exponential ($\beta = 1$) and present-biased ($\beta < 1$) preferences. Under exponential discounting, households discount future utility consistently over time, such that lifetime utility is given by

\begin{equation}
U_t=\sum_{s=t}^{T}\delta^{\,s-t}u(C_s).
\end{equation}

In contrast, present-biased households place relatively more weight on current utility than on future utility. Their preferences follow quasi-hyperbolic discounting:

\begin{equation}
U_t=u(C_t)+\beta\sum_{s=t+1}^{T}\delta^{\,s-t}u(C_s),
\end{equation}

where $\beta<1$ captures present bias. The paper focuses on naïve present-biased households, who do not fully anticipate that their future selves will also display present bias. As a result, households systematically mispredict their future consumption and saving behavior over the lifecycle.

\subsection{Estimation Strategy}

The paper estimates household preference parameters using a two-stage Method of Simulated Moments (MSM) procedure. In the first stage, the authors estimate the economic environment faced by households, including demographic profiles, borrowing limits, interest rates, and the stochastic income process. In the second stage, these estimated processes are incorporated into a lifecycle simulation model in order to estimate the preference parameters.

The first stage uses data from several U.S. datasets to estimate how household characteristics evolve over the lifecycle. In particular, the paper estimates age profiles for the number of children, dependent adults, and credit limits, as well as the income process faced by households. Labor income is modeled as a combination of a persistent AR(1) component and a transitory i.i.d. shock. The first-stage estimates are treated as fixed inputs in the second-stage estimation.

In the second stage, the model simulates household behavior over the lifecycle for different combinations of preference parameters $\theta = (\beta, \delta, \rho)$ and calculates a total of 16 moments (4 moments in 4 groups of age). The simulated moments are then compared to empirical moments calculated from the data. The preference parameters are chosen by minimizing the distance between simulated and empirical moments.

The empirical moments used for identification mainly relate to household debt and wealth accumulation patterns over the lifecycle. These include the fraction of households carrying credit card debt, the average amount of debt among indebted households, and average wealth levels for households with and without debt across different age groups. The moments are calculated for four age groups ranging from 21 to 60 years old: 21--30, 31--40, 41--50, and 51--60. The paper simulates 10,000 households in order to compute these moments.

\subsection{Data}

The empirical moments used in the paper are constructed primarily from U.S. household survey data. The main dataset is the Survey of Consumer Finances (SCF) covering the period from 1995 to 2013. The SCF is used to measure household debt, wealth, and asset holdings across the lifecycle. The analysis focuses on households whose head has a high school degree but no college degree, with additional sample restrictions including positive income, credit card ownership, and no self-employed household head.

To estimate the income process, the paper uses data from the Panel Study of Income Dynamics (PSID) between 1982 and 1991. These data are used to estimate the persistent and transitory components of labor income, as well as unemployment profiles across educational groups.

Demographic profiles, including the number of children and dependent adults over the lifecycle, are estimated using IPUMS survey data from 2001 to 2014. Additional datasets are used for specific calibrations. The Federal Reserve Economic Data (FRED) and the Bureau of Labor Statistics (BLS) provide inflation and unemployment measures, the Social Security Administration (SSA) provides survival probabilities, and the American Bankruptcy Institute (ABI) is used for the calibration of borrowing interest rates.

%re-chequear que esto este bien

\subsection{Computational Implementation}

The original project combines the use of Stata and Matlab. Stata is primarily used for data cleaning, processing, and the estimation of first-stage parameters. In particular, the authors use Stata to process the IPUMS data and estimate demographic profiles, to work with the PSID panel data and estimate components of the income process, and to clean and prepare the SCF data used to construct household debt and wealth moments.

Matlab is used for the computationally intensive components of the project. In particular, the estimation of the AR(1) income process, the variance-covariance matrix calculations, the calibration of interest rates, and the MSM estimation procedure are implemented in Matlab. Most importantly, the lifecycle simulation model itself is fully implemented manually in Matlab, including the backward induction procedure, forward simulation, grid construction, and the stochastic processes governing household decisions.

This implementation provides substantial flexibility over the structure of the model, but it also requires manually coding and managing all numerical components of the simulation. As a result, extending or modifying the model can become computationally and technically demanding.

\subsection{Main Findings}

The paper finds strong evidence in favor of present-biased preferences. The estimated present-bias parameter is $\hat{\beta}=0.530$, implying that households place substantially more weight on current utility relative to future utility. The hypothesis of exponential discounting ($\beta = 1$) is rejected by the data. The estimated long-run discount factor is $\hat{\delta}=0.989$, while the estimated coefficient of relative risk aversion is $\hat{\rho}=1.936$.

The results show that the present-biased model is able to closely match both the observed credit card borrowing behavior and the lifecycle accumulation of wealth observed in the data. In particular, the model reproduces the high prevalence of credit card debt among households while simultaneously generating positive illiquid asset holdings.

The paper also estimates a restricted version of the model imposing exponential discounting ($\beta = 1$). Under this specification, the model is still capable of reproducing the general wealth accumulation profile over the lifecycle. However, it fails to match the observed credit card borrowing patterns, substantially underpredicting the proportion of households carrying debt.

Overall, the results suggest that present bias plays an important role in explaining household borrowing and saving decisions over the lifecycle.

% todo esto doble revisar con el paper original...

\subsection{Relevance for This Thesis}

This paper represents a suitable case for computational replication due to both its economic relevance and its complex numerical implementation. The model combines lifecycle consumption and saving decisions with present-biased preferences, generating rich household behavior that depends heavily on the computational structure of the simulation. In addition, the authors provide a detailed replication package, making the paper particularly well suited for reproducibility analysis.

The paper is also relevant from a computational perspective. The original implementation relies on manually coded numerical procedures in Matlab, including backward induction, forward simulation, and the management of multiple stochastic processes and state grids. Replicating this type of model in a different computational environment provides an opportunity to evaluate how sensitive the results are to implementation choices and numerical assumptions.

This thesis reimplements the model in Python using the PyLCM package, which provides a more modular framework for defining lifecycle models. The objective is not to modify the economic structure of the original model, but rather to evaluate whether the main results can be reproduced under an alternative computational implementation. In doing so, the thesis also explores the practical challenges of implementing a complex lifecycle model within a modular computational framework.


%%%%%% NEXT SECTION %%%%%%

\section{Replication Methodology}
\label{sec:replication}

\subsection{Replication Strategy}

This thesis conducts a close computational replication of \citet{laibson2024estimating}. The objective is to reproduce the lifecycle simulation model and evaluate whether similar household debt and wealth patterns can be obtained under an alternative computational implementation.

The original model, implemented in Matlab and Stata, is reimplemented in Python using the PyLCM package. Rather than modifying the economic structure of the model, the replication focuses on translating the original lifecycle problem into a modular framework based on explicitly defined states, actions, constraints, and stochastic processes.

Due to the computational cost of the Method of Simulated Moments (MSM) estimation procedure, this thesis does not re-estimate the preference parameters of the model. Instead, the simulations are conducted using the parameter values reported in the original paper, and the evaluation focuses on comparing the simulated lifecycle moments with those reported by the authors.

\subsection{Model Implementation}

The implementation of the model required translating the original Matlab lifecycle simulation into the structure imposed by PyLCM. Although the economic structure of the model remains unchanged, implementing it in PyLCM required a very different computational structure from the original Matlab code. In the original implementation, the lifecycle problem is solved through manually coded backward induction and forward simulation routines. In contrast, PyLCM requires the model to be formulated explicitly in terms of states, actions, constraints, and stochastic processes.

\subsubsection{Functions}

One of the main differences between the original Matlab implementation and PyLCM is the overall structure of the code. In the original implementation, the focus is on explicitly coding the backward induction and forward simulation procedures. In PyLCM, instead, the model is built by defining a set of functions that describe the economic environment, while the package handles the solution process internally.

For this replication, I organized the implementation into three main groups of functions.

The first group corresponds to the regime transition functions. To understand these functions, it is important to first define what a regime means in PyLCM. A regime represents a phase of the lifecycle in which the relevant functions or parameters of the model remain unchanged. In this implementation, I defined three regimes: working life, retirement, and death. While distinguishing between working life and retirement was not strictly necessary for reproducing the baseline results of the paper, I decided to separate them in order to keep the implementation flexible for possible future extensions. The death regime differs substantially from the other regimes because households no longer make decisions or receive utility once they die. In PyLCM terminology, this corresponds to a terminal regime, meaning that households cannot transition to another regime afterwards. The regime transition functions therefore determine both how households move across regimes over the lifecycle and how state variables evolve from one period to the next.

The second group corresponds to the constraint functions. These functions define the feasible set of decisions available to households in each period. In practice, they return boolean conditions that assign a value of $-\infty$ to infeasible choices during the optimization process. One important example is the consumption constraint implied by the household budget constraint, which guarantees that consumption remains positive. More generally, these functions restrict households to economically feasible decisions, such as respecting borrowing limits or satisfying non-negative wealth conditions at the end of life.

The last group corresponds to the model functions themselves. These are the functions that define the economic structure of the model, including the utility function, the earnings process, household composition profiles, survival probabilities, and other lifecycle components required by the simulation. In this framework, the economic environment is defined directly through functions that closely resemble the mathematical structure of the model.

\subsubsection{Grids}

In the original paper, the only explicitly defined grids correspond to total liquid and illiquid asset holdings in the following period. Investment decisions are obtained implicitly from the difference between current and future asset positions. In PyLCM, however, actions and states must be defined separately.

PyLCM separates the model into state variables and action variables. State variables represent the outcomes inherited from previous decisions and shocks, while action variables represent the choices households make in each period in order to maximize utility subject to the model constraints. In this model, total liquid wealth $X_t$ and total illiquid wealth $Z_t$ are modeled as state variables, while investment decisions in liquid and illiquid assets, $I_t^X$ and $I_t^Z$, are modeled as actions. The package therefore requires explicit grids both for the state variables and for the actions available to households in each period.

The construction of the grids differs from the original implementation in several ways. In the Matlab code, the liquid asset grid varies with age because the borrowing constraint depends on the household lifecycle profile. More specifically, the minimum feasible level of liquid wealth changes over time according to age-dependent credit limits, generating a different liquid wealth grid for each age. In PyLCM, however, grids are defined at the regime level rather than separately for each age. For this reason, instead of constructing age-specific liquid wealth grids, I implemented a single liquid wealth grid together with an explicit borrowing constraint that varies over the lifecycle and restricts the feasible level of total liquid wealth in each period.

Similarly, explicit action grids for liquid and illiquid investment decisions had to be introduced, since PyLCM evaluates the optimization problem over predefined action spaces. The ranges and number of grid points were selected to balance computational feasibility with the ability to approximate the original model as closely as possible.

\subsubsection{Stochastic Processes}

This model includes two main sources of stochasticity: stochastic regime transitions and stochastic income shocks.

In the original paper, survival probabilities are incorporated directly into the backward induction problem. Households therefore make consumption and saving decisions while taking into account the possibility of dying the next periods. During the forward simulation, however, households do not actually die. Instead, after computing the simulated moments, the authors weight the moments using survival probabilities in order to avoid overrepresenting older households in the sample.

PyLCM handles this process differently through stochastic regime transitions. In this framework, households transition probabilistically to the death regime according to a Markov process using the survival probabilities estimated from the data. As a result, some households effectively transition into death during the simulation itself, removing the need to adjust the moments afterwards for survival probabilities.

The model also includes stochastic income shocks. In the original Matlab implementation, these processes are discretized manually and integrated directly into the backward induction and forward simulation routines. In PyLCM, these stochastic processes can be defined more directly by specifying the type of process and its parameters.

For the permanent component of income, I used an AR(1) process discretized through the Tauchen method, as the original paper. For the transitory component of income, it was used an i.i.d. normal shock, which captures temporary income fluctuations that do not persist over time. In both cases, PyLCM internally constructs the discretized states and transition probabilities required for the solution of the model.

\subsection{Numerical and Computational Challenges}

Implementing the model in PyLCM generated several numerical and computational challenges. Most of these issues emerged from differences in how the optimization problem is structured and solved computationally relative to the original Matlab implementation.

\subsubsection{Feasible State Space and Grid Bounds}

One of the main challenges was defining feasible lower bounds for the state grids. In the original Matlab implementation, feasible actions are implicitly restricted through the next-period state grids. In PyLCM, however, the backward induction evaluates every possible combination of states explicitly. As a result, some combinations of state variables generated optimization problems with no feasible action.

The main issue emerged at older ages, particularly near the terminal period where households are not allowed to die with negative liquid wealth. Initially, I attempted to construct the liquid wealth grid using lower bounds close to those implied by the borrowing limits in the original paper. However, for some combinations of liquid wealth, illiquid wealth, and income realizations, households were unable to satisfy simultaneously the positive consumption constraint and the no-debt-at-death condition.

For example, consider a household at age $t=90$ with income $Y_{90}=6925$, liquid wealth $X_{90}=-28,867$, and illiquid wealth $Z_{90}=0.1$. In this situation, the household must choose investment actions that eliminate negative liquid wealth before death. However, given the low level of income and the almost nonexistent illiquid wealth available for liquidation, no feasible combination of actions satisfies the budget constraint while maintaining positive consumption.

To address this problem, I modified both the lower bound of the liquid wealth grid and the minimum level of illiquid wealth. The objective was to guarantee that households would always have enough available resources to satisfy terminal constraints even under low-income realizations. In the final specification, the minimum values for liquid and illiquid wealth were set to $-5,000$ and $2,000$ respectively.

Although these adjustments resolved most infeasible cases, a very small number of households, less than 0.1\% of simulations, still encountered invalid states near the end of life during the forward simulation. One possible explanation is related to differences between the backward induction and the forward simulation processes. During backward induction, the income process is approximated using a discrete three-state Markov chain, while during forward simulation households receive shocks from the underlying continuous AR(1) process. Because the permanent shock is persistent, some households may experience sequences of negative shocks that move them into regions of the state space that were not fully represented during backward induction. In those cases, households may end up making decisions based on policy functions evaluated at nearby grid points rather than at their exact realized state, which can eventually generate infeasible terminal situations. Importantly, these cases represent only a very small fraction of simulated households and do not qualitatively affect the aggregate lifecycle patterns of the model.

\subsubsection{Action Grid Limitations}

Another important challenge was the construction of the action grids for liquid and illiquid investment decisions.

Unlike the original Matlab implementation, PyLCM requires explicit action grids over which the optimization problem is solved. The size and range of these grids, therefore, directly affect both computational feasibility and numerical accuracy.

Initially, I attempted to define wide action grids covering most of the state space. However, this quickly generated memory limitations due to the large number of state-action combinations evaluated during backward induction. Because of these computational restrictions, it was necessary to reduce the number of grid points and narrow the feasible action ranges.

The final implementation used 100 grid points for both liquid and illiquid investment decisions, with ranges of $[-50{,}000,\ 50{,}000]$ and $[-100{,}000,\ 100{,}000]$ respectively. % ACA NO ES 100& ASI ACOMODAR 

While these restrictions made the model computationally possible, they also introduced additional rigidity into household decisions. In particular, households could not adjust their portfolios as freely as in the original implementation. This limitation is likely one of the main reasons why my simulation generates larger illiquid wealth accumulation and higher debt participation relative to the original paper.

\subsubsection{Boundary Problems}

A final numerical issue emerged from the treatment of states outside the predefined grids.

If a household reaches a state value outside the grid bounds during the forward simulation, the value function is evaluated at the closest boundary point rather than at the true state value. This creates distortions because households effectively make decisions using value functions that were never computed for their actual state realizations.

To reduce this problem, I introduced additional end-of-period constraints preventing households from moving outside the feasible grid region. This ensured that the value function was always evaluated within the state space covered during backward induction.


\subsection{Evaluation Strategy}

Due to the computational cost of the MSM estimation procedure, this thesis does not re-estimate the structural preference parameters of the model. Instead, the evaluation focuses on whether the PyLCM implementation is capable of reproducing moments and lifecycle patterns similar to those reported in the original paper.

To evaluate the stability of the simulation results, I repeated the simulation multiple times using different random seeds. For each run, the model generated the simulated moments that were studied in the original paper. I then computed the mean and standard deviation of these simulated moments across repetitions.

Using these simulated distributions, I compared my results with the moments reported in the original paper by computing standardized differences between both sets of moments. This allows me to evaluate whether the observed discrepancies are consistent with stochastic simulation variability or whether they reflect more systematic differences across implementations. At the same time, some discrepancies may also be related to possible errors or misinterpretations in my own implementation of the original Matlab code, given the complexity of translating the model into a different computational framework.

\subsection{Data}

This replication does not modify the underlying data or the empirical calibration strategy of the original study. The same first-stage estimates and calibrated profiles reported in the original paper were used as inputs for the simulation model. 

The objective of this thesis is therefore not to evaluate differences generated by alternative data construction or estimation procedures, but rather to assess how the model behaves under a different computational implementation. For this reason, the second-stage simulated moments obtained in this thesis are directly compared with the moments reported in the original paper.

%%%%%%%%% next section %%%%%%%%%%%

\section{Results and Discussion}
\label{sec:results}

This section presents the main results obtained from the replication exercise and compares them with the results reported in the original paper. I focus on two specifications of the model: exponential preferences ($\beta = 1$) and naïve present-biased preferences ($\beta < 1$).

As mentioned in section \ref{sec:replication}, I do not estimate the parameters $\beta$, $\delta$, and $\rho$ given time limitations. Therefore, the objective of this section is not to estimate a new set of parameters, but rather to evaluate how closely the PyLCM implementation reproduces the original model dynamics and to analyze the sources of divergence between both implementations.

\subsection{Simulation Results and Discussions}

\subsubsection{Comparison of Simulated Moments}

As in the original paper, I computed moments related to household debt participation, average debt levels, and wealth accumulation over the lifecycle. The comparison between the moments reported in the original paper and those obtained from my simulation is presented in Tables \ref{tab:replication_exponential} and \ref{tab:replication_presentbiased} for the exponential and naïve present-biased specifications, respectively.

\begin{table}[H]
\centering
\caption{Replication Check: Exponential Model}
\label{tab:replication_exponential}
\small
\begin{tabular}{lrrrrr}
\toprule
 & Own Simulation & Paper Simulation & Std. Error & t-Statistic & p-Value \\
\midrule
\% Visa 21--30    & 0.316 & 0.309 & 0.002087 & 3.536    & 0.000 \\
\% Visa 31--40    & 0.384 & 0.287 & 0.001980 & 48.952   & 0.000 \\
\% Visa 41--50    & 0.431 & 0.299 & 0.002112 & 62.350   & 0.000 \\
\% Visa 51--60    & 0.376 & 0.257 & 0.002324 & 51.115   & 0.000 \\
mean Visa 21--30  & 0.027 & 0.044 & 0.000260 & $-$63.671  & 0.000 \\
mean Visa 31--40  & 0.025 & 0.051 & 0.000240 & $-$108.128 & 0.000 \\
mean Visa 41--50  & 0.024 & 0.059 & 0.000227 & $-$152.977 & 0.000 \\
mean Visa 51--60  & 0.021 & 0.034 & 0.000176 & $-$72.131  & 0.000 \\
wealth 21--30 $|$ debt    & 2.163 & 0.880 & 0.005770 & 222.444  & 0.000 \\
wealth 31--40 $|$ debt    & 2.938 & 0.999 & 0.022052 & 87.938   & 0.000 \\
wealth 41--50 $|$ debt    & 4.892 & 1.941 & 0.036483 & 80.890   & 0.000 \\
wealth 51--60 $|$ debt    & 8.010 & 3.811 & 0.058435 & 71.855   & 0.000 \\
wealth 21--30 $|$ no debt & 2.713 & 2.017 & 0.005576 & 124.731  & 0.000 \\
wealth 31--40 $|$ no debt & 3.093 & 2.850 & 0.014390 & 16.862   & 0.000 \\
wealth 41--50 $|$ no debt & 4.778 & 3.967 & 0.031199 & 25.987   & 0.000 \\
wealth 51--60 $|$ no debt & 7.767 & 5.358 & 0.059605 & 40.417   & 0.000 \\
\bottomrule
\end{tabular}
\vspace{0.5em}
\begin{minipage}{\linewidth}
\small This table compares the 16 simulated moments obtained from our replication of the exponential discounting model against the corresponding values reported in the original paper. Standard errors and t-statistics test equality between both sets of moments and are computed over 100 repetitions of the simulation.
\end{minipage}
\end{table}

\begin{table}[H]
\centering
\caption{Replication Check: Present-Biased Model}
\label{tab:replication_presentbiased}
\small
\begin{tabular}{lrrrrr}
\toprule
 & Own Simulation & Paper Simulation & Std. Error & t-Statistic & p-Value \\
\midrule
\% Visa 21--30    & 0.970 & 0.605 & 0.000582 & 627.572   & 0.000 \\
\% Visa 31--40    & 0.916 & 0.585 & 0.001305 & 253.564   & 0.000 \\
\% Visa 41--50    & 0.878 & 0.523 & 0.001458 & 243.126   & 0.000 \\
\% Visa 51--60    & 0.845 & 0.475 & 0.001649 & 224.196   & 0.000 \\
mean Visa 21--30  & 0.128 & 0.103 & 0.000453 & 54.429    & 0.000 \\
mean Visa 31--40  & 0.080 & 0.117 & 0.000472 & $-$77.627  & 0.000 \\
mean Visa 41--50  & 0.065 & 0.124 & 0.000346 & $-$170.188 & 0.000 \\
mean Visa 51--60  & 0.064 & 0.116 & 0.000333 & $-$155.848 & 0.000 \\
wealth 21--30 $|$ debt    & 2.340 & 0.913 & 0.003430 & 416.048   & 0.000 \\
wealth 31--40 $|$ debt    & 2.848 & 1.412 & 0.017577 & 81.725    & 0.000 \\
wealth 41--50 $|$ debt    & 5.155 & 2.640 & 0.037016 & 67.935    & 0.000 \\
wealth 51--60 $|$ debt    & 9.354 & 4.723 & 0.070570 & 65.628    & 0.000 \\
wealth 21--30 $|$ no debt & 2.821 & 2.324 & 0.026703 & 18.608    & 0.000 \\
wealth 31--40 $|$ no debt & 4.125 & 3.248 & 0.060757 & 14.427    & 0.000 \\
wealth 41--50 $|$ no debt & 6.733 & 4.633 & 0.082841 & 25.348    & 0.000 \\
wealth 51--60 $|$ no debt & 11.043 & 7.475 & 0.119768 & 29.787   & 0.000 \\
\bottomrule
\end{tabular}
\vspace{0.5em}
\begin{minipage}{\linewidth}
\small This table compares the 16 simulated moments obtained from our replication of the present-biased discounting model against the corresponding values reported in the original paper. Standard errors and t-statistics test equality between both sets of moments and are computed over 100 repetitions of the simulation.
\end{minipage}
\end{table}



In both specifications, my simulation generates a larger fraction of households holding debt across all age groups compared to the original paper. Although debt participation still decreases over the lifecycle, the differences are substantially larger in the present-biased model. At the same time, average liquid debt levels are generally lower in my implementation, with differences also bigger for the present-biased model.

One possible explanation is related to the restrictions imposed by the state and action grids. In my implementation, households face tighter limits on how much liquid wealth they can borrow in a single period and on how much wealth they can accumulate or decumulate over time. While the original paper allows households to move more freely across the state space, the PyLCM implementation requires explicit action grids with finite bounds. Because of computational limitations, these grids needed to be reduced, which decreases the accuracy of the approximation. However, the strongest restriction comes from the state grids. Households should theoretically be allowed to accumulate more debt than they do in my implementation, but if my interpretation and implementation of the model are correct, allowing for a substantially more negative liquid wealth grid generates infeasible states and numerical errors, as explained in Section \ref{sec:replication}. As a result, some households that would optimally choose higher levels of liquid borrowing are instead pushed toward larger holdings of illiquid wealth.

This mechanism also helps explain why illiquid wealth accumulation is substantially higher in my simulations for both, debt and no debt holders. Since illiquid assets generate utility benefits while also being costly to liquidate, households may use them as a substitute when liquid borrowing possibilities become too restricted. However, once households accumulate large amounts of illiquid wealth, adjusting their portfolios becomes more difficult because selling illiquid assets is costly and, given the action restrictions explained in Section \ref{sec:replication}, large adjustments cannot be made quickly. This creates additional rigidity in the model dynamics. As a result, households may simultaneously hold more illiquid wealth while also relying more heavily on liquid debt during certain stages of the lifecycle. 

The differences become considerably larger in the present-biased specification. In this case, households systematically underestimate their future consumption needs and therefore save less liquid wealth than would be optimal from the perspective of their future selves.

In the original paper, households can adjust their portfolios more flexibly over time. In my implementation, however, the restricted action grids make these adjustments more difficult. As a result, the consequences of early mistakes accumulate more strongly over the lifecycle. Once households realize that future consumption is higher than previously anticipated, they may no longer be able to adjust their asset positions sufficiently fast. This leads to higher debt participation and larger holdings of illiquid wealth relative to the original paper.


\subsubsection{Lifecycle Profiles}

The simulated lifecycle profiles reproduce several qualitative patterns observed in the original paper. Average earnings increase during early adulthood, peak around middle age, and decline toward retirement. Consumption remains substantially smoother than income over the lifecycle, indicating that households use borrowing and asset accumulation to partially smooth temporary income fluctuations.

HERE GRAPHS

The main differences emerge in the composition of household wealth. Relative to the original paper, my implementation generates larger levels of illiquid wealth accumulation for both, present-biased and exponential households, and a slower decumulation of assets toward the end of life for present-biased households.

One possible explanation is that the action restrictions introduced by the explicit grids reduce households' ability to rapidly adjust their portfolios later in life. As a result, some households accumulate asset positions that become difficult to reverse, particularly when combined with present-biased decision-making and the costly liquidation of illiquid assets. These restrictions also slow the transition from debt use to illiquid asset use as a mechanism for smoothing consumption and responding to income shocks.

An interesting difference appears in the lifecycle path of wealth accumulation under exponential preferences, where liquid wealth follows a completely different pattern from the original paper. One possible explanation is that households would prefer to decumulate illiquid wealth more rapidly later in life, but they are unable to do so because of the restrictions imposed by the action grids. Since households in my implementation also accumulate substantially more illiquid wealth than in the original paper, the limits on how much wealth can be adjusted in a single period make rapid decumulation impossible. As a result, households may accumulate additional liquid wealth as a precautionary buffer in order to smooth consumption and protect themselves against earnings shocks during later stages of the lifecycle.



\section{My Extension: Present Bias and Financial Constraints}

I decided to further explore the simulation results by looking at how behavior differs across households experiencing very low levels of consumption.

One important discussion in the behavioral economics literature is whether some behaviors commonly interpreted as present bias may instead partially reflect strong financial constraints. In particular, low-income households facing severe liquidity constraints may behave similarly to present-biased agents even when their preferences are not necessarily time-inconsistent. For example, \citet{cassidy2018poor} show that, for a group of low-income households in Pakistan, measured present bias becomes substantially smaller once financial constraints are relaxed.

Motivated by this idea, I explored whether some households in the simulation experience consumption levels that may be unrealistically low. The moments in the original paper are estimated using households without higher education, meaning that some simulated households may receive extremely negative income shocks during certain periods of the lifecycle. In some cases, these shocks could leave households consuming at levels that may not be consistent with a reasonable minimum standard of living.

To explore this, I compared simulated consumption levels with the official U.S. poverty thresholds for 2013. According to the U.S. Census Bureau, the poverty threshold was $\$15,996$ for a household without children and $\$24,421$ for a household with two children under 18. MIRAR ESTA CITA PORQUE ESTA COPY-PASTED\footnote{The official poverty definition uses money income before taxes and excludes non-cash benefits such as food assistance and housing subsidies. The thresholds do not vary geographically. Source: U.S. Census Bureau (2014).} Although these thresholds are obviously simplified measures and do not fully capture household needs, they can still provide a rough benchmark for minimum consumption requirements.

Figure XX shows that some households in the simulation consume below these thresholds during certain periods of the lifecycle, in the exponential preference specification. While these households represent a relatively small fraction of the simulated population, the result suggests that the model may allow unrealistically low levels of consumption. Since the utility specification only requires consumption to remain strictly positive, households may choose extremely small consumption values instead of borrowing more or adjusting their portfolios differently. This does not necessarily imply that the present-biased interpretation of the original paper is incorrect. However, it raises the question of whether introducing a minimum consumption floor could alter some of the simulated lifecycle behaviors. In particular, if households were required to maintain a more realistic subsistence level of consumption, some of them might need to rely more heavily on debt or liquidate assets earlier in the lifecycle.

A possible direction for future research would be to explore lifecycle models that explicitly incorporate consumption floors or baseline consumption needs. For example, \citet{noor2009hyperbolic} introduce a baseline consumption component through CARA utility. Extending the present framework in this direction could help evaluate whether some of the simulated wealth and debt dynamics remain robust once households are prevented from sustaining unrealistically low consumption levels.

HERE GRAPH

\section{Conclusion}
\label{sec:Conclusion}

The replication exercise highlights how sensitive lifecycle models can be to numerical implementation choices. Even when the economic structure, calibrated parameters, and theoretical assumptions remain unchanged, modifications in grid construction, stochastic processes, state transitions, and computational restrictions can substantially affect the simulated outcomes.

This sensitivity becomes particularly important in models with present-biased preferences, where small differences in intertemporal decisions may accumulate over the lifecycle and generate large differences in wealth dynamics.

At the same time, it is important to recognize that some of the discrepancies between my results and those reported in the original paper may also reflect limitations or mistakes in my own implementation. Translating a manually coded Matlab model into a modular framework such as PyLCM required many technical and numerical decisions regarding grids, constraints, regimes, transition functions, and stochastic processes. Therefore, differences in the interpretation of the original code, simplifications introduced during implementation, or mistakes in the specification of functions and constraints may also contribute to the observed deviations.

More generally, the project illustrates both the advantages and limitations of modular computational frameworks such as PyLCM. While the package substantially improves readability, modularity, and reproducibility, translating an existing hardcoded lifecycle model into a generalized framework also requires additional numerical decisions that can alter the behavior of the simulation in non-trivial ways.

% -------------------------------------------------------
% BIBLIOGRAPHY
% -------------------------------------------------------
\newpage
\printbibliography

% -------------------------------------------------------
% APPENDIX (optional, Roman numerals)
% -------------------------------------------------------
\newpage
\appendix
\pagenumbering{Roman}
\setcounter{page}{1}

% ... appendix content here ...

% -------------------------------------------------------
% STATEMENT OF AUTHORSHIP (last separate sheet)
% -------------------------------------------------------
\newpage
\pagenumbering{gobble}  % no page number on this sheet
\section*{Statement of Authorship}
\thispagestyle{empty}

I hereby confirm that the work presented has been performed and interpreted solely by myself except for where I explicitly identified the contrary. I assure that this work has not been presented in any other form for the fulfillment of any other degree or qualification. Ideas taken from other works in letter and in spirit are identified in every single case.

\vspace{2cm}
\noindent
Place, Date: \underline{\hspace{5cm}} \hfill Signature: \underline{\hspace{5cm}}

\end{document}